% ONINU study notes -- generated by _work/pipeline/build_tex.py from markdown.
% Compile with:  latexmk -xelatex plan.tex   (lualatex also works)
\documentclass[11pt,a4paper,parskip=half,headsepline,numbers=noenddot]{scrartcl}

% ---------------------------------------------------------------- page layout
\usepackage[a4paper,left=2cm,right=2cm,top=2.4cm,bottom=2.2cm,headheight=15pt,headsep=0.5cm]{geometry}

% ---------------------------------------------------------------- math
\usepackage{amsmath,amssymb}
\usepackage{mathtools}
\usepackage{libertinust1math}  % Libertinus math, matches the text font
\usepackage{bm}

% ---------------------------------------------------------------- fonts and language
\usepackage[no-math]{fontspec}  % no-math: keep the Type1 math fonts set up above
\IfFontExistsTF{LibertinusSerif-Regular.otf}{%
  \setmainfont{LibertinusSerif}[Extension=.otf, UprightFont=*-Regular, BoldFont=*-Bold,
    ItalicFont=*-Italic, BoldItalicFont=*-BoldItalic]
  \setsansfont{LibertinusSans}[Extension=.otf, UprightFont=*-Regular, BoldFont=*-Bold,
    ItalicFont=*-Italic, BoldItalicFont=*-Bold]
}{% fallback: TeX Gyre Pagella / Heros
  \setmainfont{texgyrepagella}[Extension=.otf, UprightFont=*-regular, BoldFont=*-bold,
    ItalicFont=*-italic, BoldItalicFont=*-bolditalic]
  \setsansfont{texgyreheros}[Extension=.otf, UprightFont=*-regular, BoldFont=*-bold,
    ItalicFont=*-italic, BoldItalicFont=*-bolditalic]
}
\setmonofont{lmmono10-regular.otf}[Scale=0.95]
\usepackage[serbian,shorthands=off]{babel}  % Serbian (Latin): Sadržaj, Slika, Tabela
\usepackage{microtype}
\setlength{\emergencystretch}{3em}
% a few Unicode characters that may be typed inside math
\usepackage{newunicodechar}
\newunicodechar{°}{\ifmmode^{\circ}\else\textdegree\fi}
\newunicodechar{≤}{\ensuremath{\leq}}
\newunicodechar{≥}{\ensuremath{\geq}}
\newunicodechar{≠}{\ensuremath{\neq}}
\newunicodechar{≈}{\ensuremath{\approx}}
\newunicodechar{∞}{\ensuremath{\infty}}
\newunicodechar{→}{\ensuremath{\rightarrow}}
\newunicodechar{⇒}{\ensuremath{\Rightarrow}}
\newunicodechar{±}{\ensuremath{\pm}}
\newunicodechar{×}{\ensuremath{\times}}
\newunicodechar{·}{\ensuremath{\cdot}}
\newunicodechar{−}{\ensuremath{-}}
\newunicodechar{ᵀ}{\ensuremath{^{\mathsf{T}}}}

% ---------------------------------------------------------------- colors, boxes
\usepackage{xcolor}
\definecolor{kljucnaC}{HTML}{1F5FA8}
\definecolor{zastoC}{HTML}{2E7D32}
\definecolor{ispitC}{HTML}{C2530A}
\definecolor{izvorC}{HTML}{6B6B6B}
\definecolor{linkC}{HTML}{1A4480}
\usepackage[most]{tcolorbox}
\tcbset{oninu/.style={enhanced, breakable, colback=#1!5, colframe=#1, colbacktitle=#1!14,
    coltitle=#1!75!black, fonttitle=\sffamily\bfseries, boxrule=0.5pt, leftrule=2.5pt,
    arc=1.5pt, left=6pt, right=6pt, top=3pt, bottom=4pt, toptitle=2pt, bottomtitle=2pt,
    before skip=8pt, after skip=10pt, parbox=false}}
% \begin{kljucna}[optional subtitle] ... \end{kljucna}   (same for zasto, ispit)
\newcommand*{\oninutitle}[2]{#1\ifx\relax#2\relax\else: #2\fi}
\newtcolorbox{kljucna}[1][]{oninu=kljucnaC, title={\oninutitle{Ključna ideja}{#1}}}
\newtcolorbox{zasto}[1][]{oninu=zastoC, title={\oninutitle{Zašto?}{#1}}}
\newtcolorbox{ispit}[1][]{oninu=ispitC, title={\oninutitle{Na ispitu}{#1}}}
% Source tag, e.g. \izvor{Khalil 13.2} -> [Khalil 13.2] in small grey type
\DeclareRobustCommand{\izvor}[1]{{\small\color{izvorC}[#1]}}

% ---------------------------------------------------------------- lists, quotes, code
\usepackage{enumitem}
\setlist{itemsep=1pt, topsep=3pt, parsep=1pt}
\providecommand{\tightlist}{}
\renewenvironment{quote}
  {\begin{tcolorbox}[enhanced, breakable, frame hidden, colback=black!4,
    borderline west={2pt}{0pt}{black!30}, sharp corners, left=8pt, before skip=6pt, after skip=8pt]\itshape}
  {\end{tcolorbox}}
\makeatletter
\def\verbatim@font{\small\ttfamily}
\makeatother

% ---------------------------------------------------------------- tables
\usepackage{longtable,booktabs,array,calc}
\newcounter{none}  % pandoc: unnumbered tables
\setlength{\LTpre}{6pt}
\setlength{\LTpost}{6pt}
\renewcommand{\arraystretch}{1.15}

% ---------------------------------------------------------------- figures
\usepackage{graphicx}
\usepackage{float}
\floatplacement{figure}{!htbp}
\usepackage[font=small, labelfont=bf, labelsep=period, format=hang, skip=4pt]{caption}
\makeatletter
\newsavebox\pandoc@box
\newcommand*\pandocbounded[1]{% scale an image to fit the text block
  \sbox\pandoc@box{#1}%
  \Gscale@div\@tempa{\textheight}{\dimexpr\ht\pandoc@box+\dp\pandoc@box\relax}%
  \Gscale@div\@tempb{\linewidth}{\wd\pandoc@box}%
  \ifdim\@tempb\p@<\@tempa\p@\let\@tempa\@tempb\fi
  \ifdim\@tempa\p@<\p@\scalebox{\@tempa}{\usebox\pandoc@box}%
  \else\usebox{\pandoc@box}\fi}
\makeatother

% ---------------------------------------------------------------- shortcuts (optional in the notes)
\providecommand{\R}{\mathbb{R}}
\providecommand{\abs}[1]{\left\lvert #1\right\rvert}
\providecommand{\norm}[1]{\left\lVert #1\right\rVert}
\providecommand{\sgn}{\operatorname{sign}}
\providecommand{\sat}{\operatorname{sat}}
\providecommand{\tr}{\operatorname{tr}}
\providecommand{\rank}{\operatorname{rank}}
\providecommand{\diag}{\operatorname{diag}}
\providecommand{\argmin}{\operatorname*{arg\,min}}
\providecommand{\argmax}{\operatorname*{arg\,max}}

% ---------------------------------------------------------------- headings, header/footer
\setkomafont{disposition}{\sffamily\bfseries\color{kljucnaC!70!black}}
\setkomafont{title}{\sffamily\bfseries\color{kljucnaC!70!black}}
\setcounter{secnumdepth}{3}
\setcounter{tocdepth}{2}
\usepackage{scrlayer-scrpage}
\clearpairofpagestyles
\ihead{\small\sffamily ONINU — plan učenja}
\ohead{\small\sffamily\pagemark}
\pagestyle{scrheadings}

% ---------------------------------------------------------------- links
\usepackage[hyperfootnotes=false]{hyperref}
\hypersetup{colorlinks=true, linkcolor=linkC, urlcolor=linkC, citecolor=linkC,
  pdftitle={ONINU — plan učenja}, pdflang={sr-Latn}, bookmarksnumbered=true, bookmarksopen=true}
\pdfstringdefDisableCommands{\def\izvor#1{}}
\usepackage{bookmark}
\urlstyle{same}

\title{ONINU --- plan učenja}
\subtitle{Pregled predmeta i redosled rada}
\author{}
\date{\today}

\begin{document}

\maketitle
\thispagestyle{plain.scrheadings}

\section{Predmet u jednoj slici}

\begin{longtable}[]{@{}
  >{\raggedright\arraybackslash}p{(\linewidth - 8\tabcolsep) * \real{0.1858}}
  >{\raggedright\arraybackslash}p{(\linewidth - 8\tabcolsep) * \real{0.2064}}
  >{\raggedright\arraybackslash}p{(\linewidth - 8\tabcolsep) * \real{0.2529}}
  >{\raggedright\arraybackslash}p{(\linewidth - 8\tabcolsep) * \real{0.2529}}
  >{\raggedright\arraybackslash}p{(\linewidth - 8\tabcolsep) * \real{0.1021}}@{}}
\caption{Tri dela, tri ispitivača, tri formata}\tabularnewline
\toprule\noalign{}
\begin{minipage}[b]{\linewidth}\raggedright
Deo
\end{minipage} & \begin{minipage}[b]{\linewidth}\raggedright
Ispitivač
\end{minipage} & \begin{minipage}[b]{\linewidth}\raggedright
Kako se polaže
\end{minipage} & \begin{minipage}[b]{\linewidth}\raggedright
Šta se predaje
\end{minipage} & \begin{minipage}[b]{\linewidth}\raggedright
Beleške
\end{minipage} \\
\midrule\noalign{}
\endfirsthead
\toprule\noalign{}
\begin{minipage}[b]{\linewidth}\raggedright
Deo
\end{minipage} & \begin{minipage}[b]{\linewidth}\raggedright
Ispitivač
\end{minipage} & \begin{minipage}[b]{\linewidth}\raggedright
Kako se polaže
\end{minipage} & \begin{minipage}[b]{\linewidth}\raggedright
Šta se predaje
\end{minipage} & \begin{minipage}[b]{\linewidth}\raggedright
Beleške
\end{minipage} \\
\midrule\noalign{}
\endhead
\bottomrule\noalign{}
\endlastfoot
Adaptivno (napredno) upravljanje & Milan Rapajić (umesto Mirne Radović) & domaći + završni projekat & MATLAB/Simulink rešenja, projekat „indirektno adaptivno upravljanje'' & \texttt{1\ -\ Adaptivno\ upravljanje} (45 str.) \\
Optimalno upravljanje & Zoran Jeličić & usmeno, fiksna lista od 12 pitanja; \textbf{uslov za izlazak: ceo domaći urađen i uredan} \izvor{P14 snimak, 51. min} & domaći = svih 12 stavki (teorija 1--5, 9--11; računski 6, 7, 8; simulacije 12) & \texttt{2\ -\ Optimalno\ upravljanje} (53 str.) \\
Nelinearno upravljanje & Željko Kanović & uradiš domaće, pokažeš ih i \textbf{braniš}; teorijska pitanja idu po domaćima & domaći 1--4 (fazni portreti, opisna funkcija, linearizacija u povratnoj sprezi, backstepping) & \texttt{3\ -\ Nelinearno\ upravljanje} (51 str.) \\
\end{longtable}

Težina ispitivača (po iskustvu kolega): Rapajić je najblaži, Jeličić i Kanović su u sredini. Kod Kanovića nije teško, ali „mora da se uloži vreme''.

\section{Redosled rada}

\begin{kljucna}
Prvo ide ono što traži najviše \textbf{rada} (domaći i projekat), a učenje napamet ide paralelno, u kraćim blokovima. Domaći se ne mogu nabubati poslednje noći; 12 pitanja mogu da se ponavljaju svaki dan po malo.
\end{kljucna}

\begin{enumerate}
\def\labelenumi{\arabic{enumi}.}
\item \textbf{Adaptivno: domaći pa projekat} (najviše kodiranja, najblaži ispitivač). Svaki domaći: pročitaj odeljak iz mape (pogl. 11, str. 44), pa uradi, pa uporedi sa referentnim rešenjima (putanje su u mapi). Projekat na kraju, prema pogl. 10 (str. 38).
\item \textbf{Nelinearno: domaći 1--4 uz pripremu odbrane.} Za svaki zadatak: postupak iz beležaka → MATLAB → pročitaj blok „Pitanja na odbrani'' (91 pitanje ukupno) i odgovori naglas bez gledanja.
\item \textbf{Optimalno: prvo domaći, pa 12 pitanja.} Domaći (sve stavke, uredno napisano) je uslov za izlazak na usmeni. Računske stavke 6, 7, 8 i 12 uradi rano (MATLAB, pogl. P6--P8, P12); teorijske stavke se pišu uz učenje, 1--2 pitanja dnevno, paralelno sa ostalim delovima. Po profesoru: P5 je „najopštiji slučaj, ko zna to zna sve'', a za P8 opšte formule za \(t_p\) i \(T\) nose mnogo. Redosled po zavisnosti: Zajednički alat → P1 → P2, P3, P4 → P5 → P9 → P10, P11 → P6, P7, P8, P12. Za svako pitanje: „Odgovor u 60 sekundi'' naglas, pa izvođenje na papiru bez gledanja, pa „Moguća potpitanja''.
\end{enumerate}

\section{Kako čitati beleške}

\begin{itemize}
\item Svaka tema ide redom: \textbf{ideja → formule → izvođenje → primer → „Na ispitu''}. Ako nešto nije jasno, pročitaj plavi okvir „Ključna ideja'' i zeleni „Zašto?''.
\item Sivi tagovi, npr. \izvor{Khalil 13.2} ili \izvor{OptControl2022 sl. 7}, upućuju na knjigu ili slajdove profesora.
\item Gde su se izvori razilazili, beleška kaže koji je tačan i zašto.
\end{itemize}

\section{Pre nego što nešto predaš: poznate zamke}

\begin{ispit}
\begin{itemize}
\item \textbf{Tvoje skripte (adaptivno):} Bode u Domaci1 koristi \texttt{log} umesto \texttt{log10}; pogrešan znak \(\hat b\) u MNK skripti; \texttt{kr}, \texttt{AA}, \texttt{CC} u Domaci5; dimenzija \texttt{CC} u Vezbe3; \texttt{rst.m} meša \(s(s-1)\) i \(s(s+1)\); \texttt{zad3} ima aliasing (sinusi do 100 Hz, a Nyquist 5 Hz).
\item \textbf{Referentno rešenje projekta:} znak jednog pola observera i nedostaje član \(B\hat w\) u observeru. Ne kopirati slepo.
\item \textbf{Česta greška (adaptivno, RST):} \(AR - SB\); tačno je \(AR + BS\).
\item \textbf{Nelinearno 2.2/2:} koristi \(G = 40/(s(s+2)(s+3))\) (zapisano sa table), ne \((s+8)\). \textbf{3.2 sistem 2:} nije striktno minimalne faze (nulta dinamika je stabilna, ali ne asimptotski).
\item \textbf{Optimalno P12:} znak pojačanja \(K\) (\(u=-Kx\) vs \(u=Kx\)) razlikuje se od P10--P11; slajd 25 odstupa od tačnog rešenja.
\end{itemize}
\end{ispit}

\section{Šta još fali (otvorena pitanja)}

\begin{itemize}
\item Tekstovi Rapajićevih domaćih: mapa u pogl. 11 je rekonstruisana iz Mirninih.
\item Kanovićeva predavanja o linearizaciji i backstepping-u: beleške se tu oslanjaju na Khalil-a; brojeve primera 14.8 i 14.9 proveriti u knjizi. Za 4.2 nema tuđeg rešenja.
\item P6 (optimalno): u tekstu je dat samo \(x(0)=1\); beleške pretpostavljaju \((1,0)\to(0,0)\). Pitati kolege šta Jeličić očekuje.
\item Transkripti predavanja (YouTube + Mirnini snimci) su u toku; kada se završe, u beleške se dodaju profesorove napomene.
\item Pisana rešenja iz referentnog git repozitorijuma (link \texttt{t.ly/\allowbreak{}Ta8OX}) još nisu preuzeta.
\end{itemize}

\section{Kontrolna lista}

\begin{itemize}
\item[$\square$]
  Adaptivno: domaći 1--6 urađeni i provereni
\item[$\square$]
  Adaptivno: projekat urađen i odbranjen
\item[$\square$]
  Nelinearno: domaći 1 (fazni portreti) + odbrana spremna
\item[$\square$]
  Nelinearno: domaći 2 (opisna funkcija) + odbrana spremna
\item[$\square$]
  Nelinearno: domaći 3 (linearizacija u povratnoj sprezi) + odbrana spremna
\item[$\square$]
  Nelinearno: domaći 4 (backstepping) + odbrana spremna
\item[$\square$]
  Optimalno: domaći (svih 12 stavki) napisan uredno i predat
\item[$\square$]
  Optimalno: P1--P5 (princip maksimuma i varijante)
\item[$\square$]
  Optimalno: P6--P8 (računski)
\item[$\square$]
  Optimalno: P9--P12 (HJB i LQR)
\item[$\square$]
  Optimalno: sva pitanja naglas, bez gledanja
\end{itemize}

\end{document}
